\section{How the New System Is Built}\label{sec:architecture}

\begin{plainwords}
We build one application, not many. It runs on our own server inside Docker containers:
a web server that handles HTTPS, the application itself, a background worker for slow
jobs like Excel exports, PostgreSQL as the single source of truth, and MinIO for files.
Every request is checked on the server: who are you, what are you allowed to do, is this
change valid in the current state, has someone else already edited this record.
\end{plainwords}

\subsection{Deployment on the VPS}

\begin{center}
\begin{tikzpicture}[
  font=\small,
  box/.style={draw, rounded corners=2pt, align=center, minimum height=9mm,
              inner sep=5pt, fill=plainbg, draw=plainframe},
  arr/.style={-{Latex[length=2.5mm]}, thick}
]
  \node[box, minimum width=6cm] (browser) {Staff browsers (any computer)};
  \node[box, below=8mm of browser, minimum width=6cm] (caddy) {\textbf{Caddy} -- reverse
    proxy, HTTPS certificates};
  \node[box, below=7mm of caddy, minimum width=6cm] (app) {\textbf{app} -- Next.js 15:
    interface + API (route handlers)};
  \node[box, below=7mm of app, minimum width=6cm] (worker) {\textbf{worker} -- same image:
    Excel exports, queue jobs (pg-boss)};
  \node[box, below=7mm of worker, minimum width=2.6cm, xshift=-1.5cm, text width=2.4cm,
        font=\footnotesize] (pg)
    {\textbf{PostgreSQL 16}\\operational data};
  \node[box, below=7mm of worker, minimum width=2.6cm, xshift=1.5cm, text width=2.4cm,
        font=\footnotesize] (minio)
    {\textbf{MinIO}\\files, labels, imports};
  \draw[arr] (browser) -- (caddy);
  \draw[arr] (caddy) -- (app);
  \draw[arr] (app) -- (worker);
  \draw[arr] (worker.west) |- (pg.north);
  \draw[arr] (worker.east) |- (minio.north);
\end{tikzpicture}
\end{center}

All five containers are defined in a single Docker Compose file, started together, and
restart on failure. Nightly jobs copy a database dump and the file store off the server.

\subsection{Inside the application: modules}

\begin{center}
\begin{tikzpicture}[
  font=\footnotesize,
  mod/.style={draw, rounded corners=2pt, align=center, fill=plainbg,
              draw=plainframe, inner sep=4pt, minimum height=8mm}
]
  \node[mod] (id) {identity\\users, roles, sessions};
  \node[mod, right=4mm of id] (cu) {customers\\master + snapshots};
  \node[mod, right=4mm of cu] (jb) {jobs\\lines, size grid, search};
  \node[mod, below=4mm of id] (st) {stock\\receipt events};
  \node[mod, right=4mm of st] (pr) {production\\stages, screens, readiness};
  \node[mod, right=4mm of pr] (aw) {artwork\\versions, approval};
  \node[mod, below=4mm of st] (sw) {swatch\\attempts, gate};
  \node[mod, right=4mm of sw] (di) {dispatch\\shipments, void};
  \node[mod, right=4mm of di] (au) {audit\\append-only logs};
  \node[mod, below=4mm of sw] (ex) {excel\\exports, staged imports};
  \node[mod, right=4mm of ex] (ig) {integrations\\outbox, DPD/Xero (later)};
\end{tikzpicture}
\end{center}

\begin{itemize}
  \item Route handlers stay thin: parse input $\rightarrow$ check permission $\rightarrow$
        call the module $\rightarrow$ serialise the reply.
  \item All business rules live inside module services -- the state machines are
        enforced in exactly one place.
  \item A static-analysis tool (dependency-cruiser) runs in CI and fails the build if a
        module reaches into another module's internals.
\end{itemize}

\subsection{The five checks every API request passes}

\begin{enumerate}
  \item \textbf{Authenticated?} -- valid server-side session.
  \item \textbf{Role permits it?} -- permission catalogue seeded in the database.
  \item \textbf{Scope and state valid?} -- department ownership, workflow transition
        allowed, field-level stripping applied.
  \item \textbf{No concurrent edit?} -- record version matches, otherwise HTTP 409 with
        the current data.
  \item \textbf{Atomic} -- entity change and its audit entry commit in one transaction.
\end{enumerate}

\subsection{Technology choices and why (plain words)}

\begin{tabularx}{\textwidth}{@{}l X@{}}
\toprule
\textbf{Choice} & \textbf{Reason} \\
\midrule
Next.js 15 (TypeScript) & One language, one code base for interface and API; huge
  ecosystem; well suited to server-side checks \\
PostgreSQL 16 & The single source of truth; strong transactions; supports the job queue
  itself \\
Drizzle ORM & Type-safe SQL with migrations kept in the repository \\
Auth.js + Argon2id + TOTP & No third-party identity service; passwords hashed with a
  memory-hard algorithm; multi-factor for key roles \\
pg-boss (no Redis) & Background jobs (exports, later integrations) queue inside
  PostgreSQL -- one less service on a small VPS \\
MinIO & S3-compatible file storage that runs on our own hardware; files stay private
  behind permission checks \\
Caddy & Automatic HTTPS certificates with almost no configuration \\
Modular monolith (no microservices) & One factory, one deployment: faster to build,
  easier to test, cheaper to run; modules keep the code clean and allow later extraction
  of integration adapters \\
\bottomrule
\end{tabularx}

\subsection*{Implementation notes for the engineering team}

\begin{itemize}
  \item \textbf{Auth.js multi-step MFA.} The Credentials provider's \texttt{authorize()}
        handles the password only -- Auth.js has no native two-step MFA. Use a staging
        flow: Step~1 validates the password and issues a short-lived,
        scope-restricted \texttt{pending\_mfa} token; Step~2 validates the TOTP against
        that token and exchanges it for the real session cookie. Users are never asked
        for a TOTP code before their password is verified, and a
        \texttt{pending\_mfa} token cannot call application endpoints.
  \item \textbf{Database security migrations.} Drizzle manages tables, columns and
        indexes -- not roles, \texttt{GRANT}/\texttt{REVOKE} or row-level security.
        Keep a separate folder of raw SQL migrations (append-only roles,
        \texttt{FORCE ROW LEVEL SECURITY} policies) executed after Drizzle migrations:
        schema migration $\rightarrow$ security migration $\rightarrow$ app deploy.
\end{itemize}